\subsection{边界承载唯一性、双指纹 Frobenius 解耦与三探针极限完备性}

正文已经分别给出：自对偶散射下绕数/Blaschke 度/真共振指数的四重同一，端点相位电流对盘内偏移总量的线性控制，\(S_{10}\times S_3\) 双指纹上的 Chebotarev 乘积律，以及素数加权二维读出的三探针最小层析。
把这三条链再作一次终端封装，可得到如下三条派生后果。

\begin{theorem}[自对偶 strictification 下的内点奇异消灭与边界承载唯一性]\label{thm:derived-selfdual-strictification-boundary-only-carrying}
\leanverified{paper\_derived\_selfdual\_strictification\_boundary\_only\_carrying}
在圆盘化--固定切片语义中，若某个散射实现满足自对偶规范并与零异常 strictification 口径相容，则一切盘内点状奇异必然消灭；准确地说，
\[
\deg(B_F)=\mathfrak d_{\mathrm{Bl}}(F)=\kappa(F)=0.
\]
因此 \(F\) 在 \(\mathbb D\) 内不再含有任何点状零点、极点或真共振指数；若系统仍保有奇异性，则其唯一容许形态是边界承载的阈值共振。并且，一旦存在可审计上界
\[
j_{\mathrm{in}}(\pi)\le M<\infty,
\]
则离临界线偏移总量满足
\[
\sum_k\delta_k\le M.
\]
\end{theorem}

\begin{proof}
由推论 \ref{cor:selfdual-implies-kappa0}，
\[
\deg(B_F)=\mathfrak d_{\mathrm{Bl}}(F)=\kappa(F)=0
\]
在自对偶散射规范下被严格强制。这消除了 \(\mathbb D\) 内全部点状零点/极点数据与真共振指数。

另一方面，备注 \ref{rem:xi-threshold-carrying-extreme-resonance} 已明确指出：一旦 \(B\equiv 1\)，盘内离散点谱共振被排除；若仍允许奇异，则它只能经由边界奇异内因子与边界测度承载，即以阈值边界共振的方式出现。

最后，由推论 \ref{cor:app-endpoint-phase-certificate} 与式后的上界段落，
\[
j_{\mathrm{in}}(\pi)\ge \sum_k\delta_k.
\]
故若 \(j_{\mathrm{in}}(\pi)\le M\)，立即得到
\[
\sum_k\delta_k\le M.
\]
证毕。
\end{proof}

\leanverified{paper\_derived\_dualfingerprint\_frobenius\_conditional\_decoupling}
\begin{theorem}[双指纹 Frobenius 事件的条件解耦]\label{thm:derived-dualfingerprint-frobenius-conditional-decoupling}
设
\[
\Gal(L_{10}L_{\mathrm{LY}}/\QQ)\cong S_{10}\times S_3.
\]
则任意来自 \(S_{10}\)-侧的共轭不变素数事件，与任意来自 \(S_3\)-侧的共轭不变素数事件，在自然密度意义下完全条件独立。特别地，若
\[
A:=\{\text{\(P_{10}\bmod p\) 含一次因子}\},
\]
\[
B_{\mathrm{split}}:=\{\text{\(g\bmod p\) 全分裂}\},
\qquad
B_{\mathrm{root}}:=\{\text{\(g\bmod p\) 至少有一根}\},
\]
则
\[
\delta(A\cap B_{\mathrm{split}})=\frac{28319}{268800},
\qquad
\delta(A\cap B_{\mathrm{root}})=\frac{28319}{67200}.
\]
又若
\[
\chi_{10}(p),\chi_{\mathrm{LY}}(p)\in\{\pm1\}
\]
为两侧判别式的二次特征，则
\[
\delta\bigl(\chi_{10}=\varepsilon_1,\chi_{\mathrm{LY}}=\varepsilon_2\bigr)=\frac14
\qquad
(\varepsilon_1,\varepsilon_2\in\{\pm1\}).
\]
因此，任一侧的二次符号层都不能改善对另一侧分裂/有根事件的预测：相应条件密度恒等于其无条件密度。
\end{theorem}

\begin{proof}
由定理 \ref{thm:xi-time-part62ea-dualfingerprint-classfunction-prime-average-factorization} 与定理 \ref{thm:xi-time-part62ea-dualfingerprint-boolean-strong-product-law}，凡分别只依赖于 \(S_{10}\)-分量和 \(S_3\)-分量的共轭不变事件都满足精确乘法密度律
\[
\delta(E\cap F)=\delta(E)\delta(F).
\]
将 \(E=A\) 与
\[
F=B_{\mathrm{split}},\qquad F=B_{\mathrm{root}}
\]
代入，再结合推论 \ref{cor:killo-leyang-joint-density-explicit}，即得
\[
\delta(A\cap B_{\mathrm{split}})=\frac{28319}{268800},
\qquad
\delta(A\cap B_{\mathrm{root}})=\frac{28319}{67200}.
\]

另一方面，由定理 \ref{thm:killo-leyang-quadratic-character-decoupling}，
\[
\bigl(\chi_{10}(p),\chi_{\mathrm{LY}}(p)\bigr)
\]
在 \(\{\pm1\}^2\) 上等分布，因此每一对符号的密度均为 \(1/4\)。

于是对任一 \(S_{10}\)-侧事件 \(E\) 与任一 \(S_3\)-侧事件 \(F\)，都有
\[
\delta(E\mid F)=\delta(E),
\qquad
\delta(F\mid E)=\delta(F),
\]
特别是二次符号层对另一侧分裂型事件没有任何条件增益。证毕。
\end{proof}

\begin{corollary}[双探针永久欠定，而三探针达到极限完备]\label{cor:derived-three-probe-minimal-completeness}
\leanverified{paper\_derived\_three\_probe\_minimal\_completeness}
在椭圆极限指纹的高斯读出口径下，任何仅依赖两条线性探针的审计都在结构上永久欠定；相反，三条探针方向的方差恰为恢复极限二次型
\[
M_\infty
\]
的最小完备数据。特别地，若两探针 \(a,b\) 满足
\[
a^\top\Sigma_\infty b=0,
\]
则其联合读出渐近独立，且四个符号象限的概率同时趋于 \(1/4\)；故二探针符号统计在正交化后完全失去混合项信息，不能辨识椭圆极限指纹。三探针完备性因此不是经验事实，而是最小识别维数。
\end{corollary}

\begin{proof}
由定理 \ref{thm:xi-time-part61aca-primeweighted-sigma-orthogonal-independence}，
\[
a^\top\Sigma_\infty b=0
\iff
X_T(a)\ \text{与}\ X_T(b)\ \text{渐近独立},
\]
并且在该条件下四个符号象限的极限概率全为 \(1/4\)。这说明：一旦把两探针正交化，符号层的混合信息完全被抹平。

另一方面，定理 \ref{thm:xi-time-part61aca-primeweighted-three-probe-minimal-tomography} 已证明：三条方向方差可以唯一恢复
\[
M_\infty=\sigma^2\Sigma_\infty,
\]
而任意只含两个方向的方差数据都不足以唯一确定该正定对称矩阵。

两者合并即得：二探针制度的缺陷来自识别维数本身不足，而三探针正是第一次达到极限完备的阈值。证毕。
\end{proof}

\endinput
